% ==============================================================
\section{Incomplete Depreciation}
\label{sec:ac-depreciation}
% ==============================================================

We now relax the assumption of full depreciation. When capital persists across periods ($\delta < 1$), the distinction between gross and net investment becomes important, and the dynamics of Tobin's $Q$ acquire a richer structure. The CES transformation technology continues to govern adjustment costs, but now applies to gross investment rather than the change in the capital stock.

% --------------------------------------------------------------
\subsection{Setup with Persistent Capital}
\label{subsec:ac-dep-setup}
% --------------------------------------------------------------

\paragraph{Capital Accumulation.}
Capital depreciates at rate $\delta \in (0, 1]$ and is augmented by gross investment $I$:
\begin{equation}
k' = (1 - \delta) k + I.
\label{eq:ac-dep-accum}
\end{equation}

Net investment is the change in the capital stock:
\begin{equation}
k' - k = I - \delta k.
\label{eq:ac-dep-net}
\end{equation}

When $\delta = 1$, we recover the full depreciation case: $k' = I$.

\paragraph{Technology.}
Production is linear in capital:
\begin{equation}
Y = A k.
\label{eq:ac-dep-production}
\end{equation}

The CES transformation governs the split of output into consumption and gross investment:
\begin{equation}
\Hfun(C, I) = \left( \omega_C^{1-\psi} C^{\psi} + \omega_I^{1-\psi} I^{\psi} \right)^{1/\psi} = Y.
\label{eq:ac-dep-transformation}
\end{equation}

\paragraph{Timing.}
At the beginning of period $t$, the household owns capital $k_t$. Production occurs, yielding output $Y_t = A k_t$. The household then chooses consumption $C_t$ and gross investment $I_t$ subject to the transformation constraint. Capital for next period is $k_{t+1} = (1-\delta) k_t + I_t$.

% --------------------------------------------------------------
\subsection{The Two-Period Problem}
\label{subsec:ac-dep-two-period}
% --------------------------------------------------------------

\paragraph{Setup.}
Consider $t \in \{0, 1\}$. The household enters period 0 with capital $k_0$, chooses $(C_0, I_0)$, and then consumes all output in period 1:
\begin{equation}
\max_{C_0, I_0} \left\{ u(C_0) + \beta \, u(C_1) \right\}
\label{eq:ac-dep-objective}
\end{equation}
subject to:
\begin{align}
\Hfun(C_0, I_0) &= A k_0, \label{eq:ac-dep-const0} \\
k_1 &= (1-\delta) k_0 + I_0, \label{eq:ac-dep-accum0} \\
C_1 &= A k_1 = A \left[ (1-\delta) k_0 + I_0 \right]. \label{eq:ac-dep-const1}
\end{align}

\paragraph{The Key Difference.}
With incomplete depreciation, future consumption depends on \emph{both} current investment $I_0$ and the undepreciated capital $(1-\delta) k_0$:
\begin{equation}
C_1 = A (1-\delta) k_0 + A I_0.
\label{eq:ac-dep-C1}
\end{equation}

The term $A(1-\delta) k_0$ is ``locked in''---it does not depend on the household's choice. Only the term $A I_0$ is affected by the investment decision.

\paragraph{First-Order Condition.}
The Lagrangian is:
\begin{equation}
\mathcal{L} = u(C_0) + \beta \, u(A[(1-\delta)k_0 + I_0]) + \lambda \left[ A k_0 - \Hfun(C_0, I_0) \right].
\label{eq:ac-dep-lagrangian}
\end{equation}

First-order conditions:
\begin{align}
C_0: & \quad u'(C_0) = \lambda \frac{\partial \Hfun}{\partial C_0}, \label{eq:ac-dep-foc-C} \\
I_0: & \quad \beta A \, u'(C_1) = \lambda \frac{\partial \Hfun}{\partial I_0}. \label{eq:ac-dep-foc-I}
\end{align}

Taking the ratio and using the definition of Tobin's $Q$:
\begin{equation}
\frac{u'(C_0)}{\beta A \, u'(C_1)} = \frac{\partial \Hfun / \partial C_0}{\partial \Hfun / \partial I_0} = \frac{1}{Q}.
\label{eq:ac-dep-ratio}
\end{equation}

The Euler equation is:
\begin{equation}
\boxed{u'(C_0) = \beta \frac{A}{Q} u'(C_1).}
\label{eq:ac-dep-euler}
\end{equation}

This has the same form as the full depreciation case! The Euler equation relates marginal utilities through the return $A/Q$, regardless of $\delta$.

\begin{calloutnote}[Why Doesn't $\delta$ Appear in the Euler Equation?]
The depreciation rate $\delta$ does not appear in the Euler equation because it affects only the \emph{level} of future consumption, not the \emph{marginal} trade-off. The household's decision at the margin is: ``Should I invest one more unit today?'' The return on that marginal unit is $A/Q$, independent of how much capital survives from the past.

However, $\delta$ does affect the \emph{solution} through the budget constraint: higher $\delta$ means less future consumption for given $I_0$, which changes the equilibrium allocation and hence $Q$.
\end{calloutnote}

% --------------------------------------------------------------
\subsection{The Finite-Horizon Problem}
\label{subsec:ac-dep-T-period}
% --------------------------------------------------------------

\paragraph{The Bellman Equation.}
For $t \in \{0, 1, \ldots, T\}$, define $V_t(k)$ as the value of entering period $t$ with capital $k$:
\begin{equation}
V_t(k) = \max_{I} \left\{ u(C) + \beta V_{t+1}((1-\delta)k + I) \right\} \quad \text{s.t.} \quad \Hfun(C, I) = A k,
\label{eq:ac-dep-bellman}
\end{equation}
with terminal condition $V_T(k) = u(Ak)$.

\paragraph{First-Order Condition.}
\begin{equation}
-u'(C) \frac{1}{Q} + \beta V_{t+1}'(k') = 0,
\label{eq:ac-dep-foc-bellman}
\end{equation}
where $k' = (1-\delta)k + I$.

\paragraph{Envelope Condition.}
Differentiating the value function:
\begin{equation}
V_t'(k) = u'(C) \frac{A}{Q} + \beta (1-\delta) V_{t+1}'(k').
\label{eq:ac-dep-envelope}
\end{equation}

The envelope condition now has two terms:
\begin{itemize}[nosep]
\item $u'(C) \cdot A/Q$: The marginal value of capital through current production (as before).
\item $\beta (1-\delta) V_{t+1}'(k')$: The marginal value of capital that survives to next period.
\end{itemize}

\paragraph{The Euler Equation.}
Combining the FOC and envelope condition (shifted forward):
\begin{equation}
\frac{u'(C_t)}{Q_t} = \beta V_{t+1}'(k_{t+1}) = \beta \left[ u'(C_{t+1}) \frac{A}{Q_{t+1}} + \beta (1-\delta) V_{t+2}'(k_{t+2}) \right].
\label{eq:ac-dep-euler-step}
\end{equation}

Iterating forward and using the FOC repeatedly:
\begin{equation}
\frac{u'(C_t)}{Q_t} = \beta u'(C_{t+1}) \frac{A}{Q_{t+1}} + \beta^2 (1-\delta) u'(C_{t+2}) \frac{A}{Q_{t+2}} + \cdots
\label{eq:ac-dep-euler-iterated}
\end{equation}

This can be written more compactly. Define the \emph{user cost of capital}:
\begin{equation}
\boxed{u'(C_t) = \beta \left[ \frac{A}{Q_{t+1}} + (1-\delta) \frac{Q_t}{Q_{t+1}} \right] Q_t \, u'(C_{t+1}).}
\label{eq:ac-dep-euler-full}
\end{equation}

\paragraph{Interpretation: The Return on Capital.}
Define the gross return on capital:
\begin{equation}
R_{t+1} \equiv \frac{A + (1-\delta) Q_{t+1}}{Q_t}.
\label{eq:ac-dep-return}
\end{equation}

The Euler equation becomes:
\begin{equation}
u'(C_t) = \beta R_{t+1} \, u'(C_{t+1}).
\label{eq:ac-dep-euler-standard}
\end{equation}

\begin{calloutimportant}[The Return with Incomplete Depreciation]
The return $R_{t+1} = (A + (1-\delta) Q_{t+1}) / Q_t$ has two components:
\begin{enumerate}[nosep]
\item \textbf{Dividend yield}: $A / Q_t$ is the marginal product per unit of capital value.
\item \textbf{Capital gain}: $(1-\delta) Q_{t+1} / Q_t$ is the value of undepreciated capital, adjusted for changes in $Q$.
\end{enumerate}

Compare with full depreciation where $R_{t+1} = A Q_t / Q_{t+1}$. With $\delta < 1$, the ``resale value'' of capital $(1-\delta) Q_{t+1}$ enters the return.

In steady state where $Q_t = Q_{t+1} = Q$:
\[
R = \frac{A + (1-\delta) Q}{Q} = \frac{A}{Q} + (1-\delta).
\]
The return equals the dividend yield plus the survival rate.
\end{calloutimportant}

% --------------------------------------------------------------
\subsection{Steady State}
\label{subsec:ac-dep-steady-state}
% --------------------------------------------------------------

Consider constant productivity $A$ and an infinite horizon (or large $T$). In steady state, all variables are constant: $C_t = C$, $I_t = I$, $k_t = k$, $Q_t = Q$.

\paragraph{Steady-State Capital Accumulation.}
From $k' = (1-\delta)k + I$:
\begin{equation}
k = (1-\delta)k + I \quad \Longrightarrow \quad I = \delta k.
\label{eq:ac-dep-ss-accum}
\end{equation}

In steady state, gross investment exactly replaces depreciated capital. Net investment is zero.

\paragraph{Steady-State Euler Equation.}
From \cref{eq:ac-dep-euler-standard} with $R_{t+1} = R$:
\begin{equation}
1 = \beta R = \beta \frac{A + (1-\delta) Q}{Q}.
\label{eq:ac-dep-ss-euler}
\end{equation}

Solving for $Q$:
\begin{equation}
Q = \frac{\beta A}{1 - \beta(1-\delta)}.
\label{eq:ac-dep-Q-ss}
\end{equation}

\paragraph{Special Cases.}

\medskip
\noindent
\textit{Full depreciation ($\delta = 1$):}
\begin{equation}
Q^{ss} = \frac{\beta A}{1 - 0} = \beta A.
\label{eq:ac-dep-Q-ss-full}
\end{equation}
This matches our earlier result.

\medskip
\noindent
\textit{No depreciation ($\delta = 0$):}
\begin{equation}
Q^{ss} = \frac{\beta A}{1 - \beta} = \frac{A}{1/\beta - 1} = \frac{A}{r},
\label{eq:ac-dep-Q-ss-none}
\end{equation}
where $r = 1/\beta - 1$ is the discount rate. This is the standard asset pricing formula: the price equals the perpetuity value of dividends.

\medskip
\noindent
\textit{General $\delta \in (0,1)$:}
\begin{equation}
Q^{ss} = \frac{A}{1/\beta - (1-\delta)} = \frac{A}{r + \delta}.
\label{eq:ac-dep-Q-ss-general}
\end{equation}

The denominator $r + \delta$ is the \emph{user cost of capital}: the sum of the interest cost and the depreciation cost.

\begin{calloutnote}[The User Cost of Capital]
The steady-state condition $Q = A / (r + \delta)$ can be rewritten as:
\[
(r + \delta) Q = A.
\]
The left-hand side is the \emph{user cost}---the cost of using one unit of capital for one period:
\begin{itemize}[nosep]
\item $r Q$: Interest cost (opportunity cost of funds tied up in capital).
\item $\delta Q$: Depreciation cost (value of capital lost to wear).
\end{itemize}
In equilibrium, the user cost equals the marginal product $A$. This is the classic Jorgensonian result, extended to include adjustment costs through $Q$.
\end{calloutnote}

\paragraph{Steady-State Investment Rate.}
From $I = \delta k$ and $Y = Ak$:
\begin{equation}
\frac{I}{Y} = \frac{\delta k}{A k} = \frac{\delta}{A}.
\label{eq:ac-dep-inv-rate}
\end{equation}

The investment-output ratio is $\delta / A$. Higher depreciation requires more investment to maintain the capital stock.

\paragraph{Steady-State Consumption.}
From the CES constraint:
\begin{equation}
\Hfun(C, I) = Y \quad \Longrightarrow \quad \Hfun(C, \delta k) = A k.
\label{eq:ac-dep-ss-ces}
\end{equation}

Given $Q$, we can solve for the consumption-investment ratio and hence for $C$ and $k$.

% --------------------------------------------------------------
\subsection{Transitional Dynamics}
\label{subsec:ac-dep-dynamics}
% --------------------------------------------------------------

Away from steady state, the system exhibits transitional dynamics governed by two state-like variables: the capital stock $k_t$ and (implicitly) the shadow value $V_t'(k_t)$.

\paragraph{The Dynamic System.}
The economy is characterized by:
\begin{align}
k_{t+1} &= (1-\delta) k_t + I_t, \label{eq:ac-dep-dyn-k} \\
\frac{u'(C_t)}{Q_t} &= \beta \left[ \frac{A}{Q_{t+1}} u'(C_{t+1}) + (1-\delta) \frac{u'(C_{t+1})}{Q_{t+1}} \right], \label{eq:ac-dep-dyn-euler} \\
\Hfun(C_t, I_t) &= A k_t. \label{eq:ac-dep-dyn-ces}
\end{align}

Given $k_t$, the CES constraint \cref{eq:ac-dep-dyn-ces} defines a relationship between $C_t$ and $I_t$. The Euler equation \cref{eq:ac-dep-dyn-euler} pins down the intertemporal allocation. Capital accumulation \cref{eq:ac-dep-dyn-k} determines $k_{t+1}$.

\paragraph{Phase Diagram Intuition.}
In the $(k, C)$ plane:
\begin{itemize}[nosep]
\item The $\dot{k} = 0$ locus (where $k_{t+1} = k_t$) requires $I = \delta k$.
\item The $\dot{C} = 0$ locus (where $C_{t+1} = C_t$) requires $\beta R = 1$, which pins down $Q$ and hence the $C/I$ ratio.
\end{itemize}

The steady state is at the intersection. The saddle-path dynamics are qualitatively similar to the standard Ramsey model, but the adjustment cost (finite $\eta$) affects the speed of convergence.

\paragraph{Effect of Adjustment Costs on Dynamics.}
Lower transformation elasticity $\eta$ (higher adjustment costs):
\begin{itemize}[nosep]
\item Slows convergence to steady state.
\item Increases the volatility of $Q$ along the transition path.
\item Reduces the response of investment to productivity shocks.
\end{itemize}

In the limit $\eta \to \infty$ (no adjustment costs), the economy jumps immediately to steady state. In the limit $\eta \to 0$ (Leontief), the consumption-investment ratio is fixed and adjustment is impossible.

% --------------------------------------------------------------
\subsection{Gross vs. Net Investment}
\label{subsec:ac-dep-gross-net}
% --------------------------------------------------------------

With incomplete depreciation, the distinction between gross and net investment matters for measurement and interpretation.

\paragraph{Definitions.}
\begin{align}
\text{Gross investment:} & \quad I_t, \\
\text{Net investment:} & \quad I_t^{\text{net}} = I_t - \delta k_t = k_{t+1} - k_t, \\
\text{Depreciation:} & \quad D_t = \delta k_t.
\end{align}

\paragraph{The CES Constraint in Terms of Net Investment.}
Substituting $I = I^{\text{net}} + \delta k$ into the transformation:
\begin{equation}
\Hfun(C, I^{\text{net}} + \delta k) = A k.
\label{eq:ac-dep-ces-net}
\end{equation}

This shows that adjustment costs apply to \emph{gross} investment, not net investment. Even maintaining the capital stock ($I^{\text{net}} = 0$) requires resources $I = \delta k$, and these resources are subject to the transformation friction.

\paragraph{Interpretation.}
The CES transformation can be viewed as capturing:
\begin{itemize}[nosep]
\item \textbf{Installation costs}: Costs of putting new capital in place.
\item \textbf{Maintenance costs}: Costs of replacing depreciated capital.
\item \textbf{Sectoral reallocation}: Costs of shifting resources between consumption and investment goods sectors.
\end{itemize}

All three apply to gross investment flows, not just to changes in the capital stock.

\begin{calloutnote}[Why Gross Investment?]
The assumption that adjustment costs apply to gross investment (rather than net) is standard in the literature. The economic rationale is that the frictions---installation, learning, organizational change---occur whenever new capital goods are produced and installed, regardless of whether they replace depreciated capital or expand the stock.

An alternative specification would have adjustment costs on net investment: $\Hfun(C, k' - (1-\delta)k) = Y$. This would imply that replacement investment is frictionless, which may be appropriate in some contexts (e.g., identical replacement of worn-out machines) but not others (e.g., technological upgrading).
\end{calloutnote}

% --------------------------------------------------------------
\subsection{Comparison: Full vs. Incomplete Depreciation}
\label{subsec:ac-dep-comparison}
% --------------------------------------------------------------

\begin{table}[H]
\centering
\caption{Adjustment Costs with Full vs. Incomplete Depreciation}
\label{tab:ac-dep-comparison}
\renewcommand{\arraystretch}{1.6}
\begin{tabular}{@{}lcc@{}}
\toprule
\textbf{Object} & \textbf{Full ($\delta = 1$)} & \textbf{Incomplete ($\delta < 1$)} \\
\midrule
Capital accumulation & $k' = I$ & $k' = (1-\delta)k + I$ \\[4pt]
Return on capital & $R = \dfrac{A \, Q_t}{Q_{t+1}}$ & $R = \dfrac{A + (1-\delta) Q_{t+1}}{Q_t}$ \\[10pt]
Steady-state $Q$ & $\beta A$ & $\dfrac{\beta A}{1 - \beta(1-\delta)} = \dfrac{A}{r + \delta}$ \\[10pt]
Steady-state $I/Y$ & Determined by $C/I$ ratio & $\delta / A$ \\[4pt]
Net investment (SS) & $I$ & $0$ \\[4pt]
\bottomrule
\end{tabular}
\end{table}

% --------------------------------------------------------------
\subsection{Summary}
\label{subsec:ac-dep-summary}
% --------------------------------------------------------------

\begin{table}[H]
\centering
\caption{Incomplete Depreciation with CES Adjustment Costs}
\label{tab:ac-dep-summary}
\renewcommand{\arraystretch}{1.6}
\begin{tabular}{@{}ll@{}}
\toprule
\textbf{Object} & \textbf{Formula} \\
\midrule
Capital accumulation & $k' = (1-\delta) k + I$ \\[6pt]
Return on capital & $R_{t+1} = \dfrac{A + (1-\delta) Q_{t+1}}{Q_t}$ \\[10pt]
Euler equation & $u'(C_t) = \beta R_{t+1} \, u'(C_{t+1})$ \\[6pt]
Steady-state $Q$ & $Q^{ss} = \dfrac{A}{r + \delta}$ where $r = 1/\beta - 1$ \\[10pt]
User cost of capital & $(r + \delta) Q = A$ \\[6pt]
Steady-state investment & $I^{ss} = \delta k^{ss}$ (replacement only) \\[6pt]
\bottomrule
\end{tabular}
\end{table}

\begin{callouttip}[Key Insights]
Incomplete depreciation introduces:
\begin{itemize}[nosep]
\item \textbf{Capital gains in returns}: The return includes the resale value $(1-\delta) Q_{t+1}$ of undepreciated capital.
\item \textbf{User cost of capital}: Steady-state $Q = A/(r + \delta)$ reflects both interest and depreciation costs.
\item \textbf{Gross vs. net investment}: Adjustment costs apply to gross investment; even replacement investment is costly.
\item \textbf{Richer dynamics}: The state variable $k_t$ evolves gradually, creating transitional dynamics even with constant productivity.
\item \textbf{Steady-state investment}: In steady state, gross investment equals depreciation ($I = \delta k$) and net investment is zero.
\end{itemize}
The CES transformation continues to govern the consumption-investment trade-off, with Tobin's $Q$ summarizing the shadow cost of capital accumulation.
\end{callouttip}
